%%%%%%%%%%%%%%%%%%%%%%%%%%%%%%%%%%%%%%%%%%%%%%%%%%%%%%%%%%
% # Title page
%%%%%%%%%%%%%%%%%%%%%%%%%%%%%%%%%%%%%%%%%%%%%%%%%%%%%%%%%%

\pagenumbering{gobble}
\centering
\vspace{3cm}

\doublespacing
THE PROCESSING AND ACCEPTABILITY OF GAPPED VS. RESUMPTIVE \\
RELATIVE CLAUSES IN FIRST AND SECOND LANGUAGE ENGLISH

\vspace{1.5cm}
\singlespacing
A DISSERTATION SUBMITTED TO THE GRADUATE DIVISION OF \\
THE UNIVERSITY OF HAWAI`I AT MĀNOA IN PARTIAL FULFILLMENT OF \\
THE REQUIREMENTS FOR THE DEGREE OF

\vspace{1cm}
\doublespacing
DOCTOR OF PHILOSOPHY \\
IN \\
SECOND LANGUAGE STUDIES

\vspace{1.5cm}
\singlespacing
AUGUST 2024

\vspace{1.5cm}
\doublespacing
By \\
Fred Zenker

\vspace{1.5 cm}
\singlespacing
Dissertation Committee: \\
Bonnie D. Schwartz, Chairperson \\
Kamil Deen \\
Li `Julie' Jiang \\
William O'Grady \\
Amy Schafer

\vspace*{\fill}
\flushleft
Keywords: relative clause, resumption, resumptive pronoun, adult second language acquisition, adult second language processing

\newpage

\pagenumbering{roman}
\setcounter{page}{2}
\vspace*{\fill}
\begin{centering}

Copyright \textcopyright \space 2024

\vspace{0.5 cm}

Fred Zenker

All Rights Reserved

\end{centering}

%%%%%%%%%%%%%%%%%%%%%%%%%%%%%%%%%%%%%%%%%%%%%%%%%%%%%%%%%%
% # Acknowledgments
%%%%%%%%%%%%%%%%%%%%%%%%%%%%%%%%%%%%%%%%%%%%%%%%%%%%%%%%%%

\chapter*{Acknowledgments}
\addcontentsline{toc}{chapter}{Acknowledgments}
\setlength{\parindent}{2.6em}
\setstretch{1.5}

Looking back on my PhD journey, I am deeply touched by all help I have received from others along they way. I would like to begin by thanking the members of my dissertation committee at the University of Hawai`i at Mānoa (UHM)---Kamil~Deen, Li~`Julie'~Jiang, William~O'Grady, Amy~Schafer, and Bonnie~D.~Schwartz---for their constant guidance. I am incredibly lucky to have completed this project under the direction of such a stellar group of researchers. First and foremost, I would like to express my gratitude to my dissertation supervisor, Bonnie~D.~Schwartz, for believing in me, for spending countless hours commenting on my work, for introducing me to her colleagues around the world, and for always challenging me to be a better researcher and scholar. I admire Bonnie for her brilliance, her apparent tirelessness, and her excellent sense of humor. There is no way I could have accomplished as much as I did at UHM without her expert tutelage. Thank you, Bonnie, for taking me under your wing for all these years and for being such a caring, generous, and understanding `intellectual parent' to me.

Kamil~Deen has offered me generous encouragement and support since I was a new arrival on campus. Thank you, Kamil, for your kind words after my conference presentations and for helping me prepare for job interviews. The `Passives' project we did in your LING 750 course remains one of my all-time favorite experiences collaborating with others.

I am thankful to Li~`Julie'~Jiang for sharing her vast knowledge of Chinese syntax with me and for offering helpful suggestions about how I could improve my Mandarin stimuli. Her detailed and timely comments on my drafts were very much appreciated. Thank you, Julie, for being so quick to offer me assistance when I needed it and for always being a friendly face around Moore~Hall.

William~O'Grady's feedback has had a profound effect on every stage of this project. I appreciate him for showing me how important it is to keep an open mind and to consider problems from a variety of perspectives. William is a very effective communicator and a constant source of inspiration for me as a try to become a better speaker and writer. Thank you, William, for taking an interest in me and for making me feel like I have something of value to contribute to the field.

Last but certainly not least, I am grateful to Amy~Schafer for teaching my about psycholinguistics and statistics. Amy is probably the committee member I consulted with most after Bonnie, and I deeply enjoyed our many office hours meetings together---as well as our occasional run-ins on the campus shuttle and at the Imin Center lānai. Thank you, Amy, for answering my many questions about data analysis and for helping me stay positive about my progress. I learned a great deal about how to be a better and more organized teacher from your excellent example.

I am also thankful to all the other instructors and mentors who have contributed to my academic and professional development, both at UHM and at other institutions. A few individuals who deserve special mention for their support of this project include Shin~Fukuda, Theres~Grüter, Daniel~Isbell, Kristopher~Kyle, Hajime~Ono, Kent~Sakoda, Jeannette~Schaeffer, Rex~A.~Sprouse, Yi-ching~Su, and Shigeo~Tonoike.

I am deeply grateful to my research assistant, Nicholas~Coney, for his help with transcription and coding. Thank you, Nick, for your flexibility and for your many hours of diligent work. I wish you the very best in your career and everything else that life has in store for you.

For help of various kinds, ranging from translation assistance and participant recruitment to moral support and statistics advice, I would like to express my sincere gratitude to the following friends and colleagues: Ivan~Bondoc, Amber~Camp, Yu-Tzu~Chang, Yunchuan~Chen, Ann~Tai~Choe, In~Ji~Sera~Chun, Daniel~Holden, Yao~Huang, Haerim~Hwang, Hyeyoung~Jung, Hyunwoo~Kim, Jieun~Kim, Jihyun~Kim, Kitaek~Kim, Rickey~Larkin, Eunsun~Lee, Wenyi~Ling, Akari~Ohba, Leeseul~Park, Anupama~Reddy, Gyu-ho Shin, Nozomi~Tanaka, Jue~Wang, Po-Jun~Yu, Mayuko~Yusa, Kim~Zenker, Jing~`Crystal'~Zhong, and Yanxin~`Alice'~Zhu. For the sheer number of hours they spent assisting me with this research, special thanks are due to Yu-Tzu~Chang, Haerim~Hwang, Kitaek~Kim, Kim~Zenker, and Jing~`Crystal'~Zhong. I feel so blessed to have all you amazing people in my life.

I am likewise indebted to all the members of the labs and research groups I have been involved in at UHM, especially the Language Analysis and Experimental Labs, the Language Acquisition Research Group, and the Quantitative Research Discussion Group. I would also like to thank Ivy~Sichel, Maziar~Toosarvandani, and Matt~Wagers at the University of California, Santa Cruz for allowing me to participate in their R/Lab meetings back in Spring 2021.

This project has received generous financial support from a number of different sources including, in order from largest to smallest contribution, the Bilinski Foundation; the National Science Foundation; the journal \textit{Language Learning}; the Dean of Languages, Linguistics, and Literature Special Fund at UHM; the Graduate Student Association Grants and Awards Program at UHM; and the Boston University Conference on Language Development. I offer my heartfelt thanks to the decision makers at each of these organizations for investing in my research. I would like to thank my UHM grant managers John Kawahara and Daniel Sato as well for their assistance with filing applications and disbursing funds.

I have been blown away by how kind and welcoming the UHM community is since the very first day I arrived on campus, and there are far more people I would like to thank than I could possibly include here. However, these acknowledgements would simply not be complete if I did not acknowledge Priscilla~Faucette, Kenton~Harsch, Jennifer~Kanda, Justin~Kanda, and Joel~Weaver for all the help they have given me over the years. Thank you all so much for looking out for me and making me feel at home in Hawai`i.

It is also due to emotional support from my family that I was able to complete this project. My deepest thanks go to my wife, Kim~Zenker, for her love, patience, and encouragement---you are the best life partner that I could ever have asked for, and I love you with all of my heart. I dedicate this dissertation to my parents, Amy~Marshall and Tim~Zenker, who raised me to believe that I can reach my goals if I put my mind to it.

Finally, I would like to thank all my participants for the time and effort they spent on my experiments. I truly could not have done it without you. 

Although I have received valuable feedback from others at every stage of the project, I am solely responsible for any errors that remain in this dissertation.

%%%%%%%%%%%%%%%%%%%%%%%%%%%%%%%%%%%%%%%%%%%%%%%%%%%%%%%%%%
% # Abstract
%%%%%%%%%%%%%%%%%%%%%%%%%%%%%%%%%%%%%%%%%%%%%%%%%%%%%%%%%%

\chapter*{Abstract}
\addcontentsline{toc}{chapter}{Abstract}

This dissertation investigates the interplay between the implicit knowledge that learners have of a nonnative language and their processing of that language, examining two types of relative clauses (RCs) in English: gapped RCs (e.g.,~\textit{the man that they hired}) and resumptive RCs (e.g.,~*\textit{the man that they hired him}). Resumptive RCs are considered ungrammatical in English, and yet they are not uncommon in nonnative speech/writing, even from learners whose first language (L1) also disallows them. Phenomena of this type represent a very interesting area of second language (L2) research because they cannot be straightforwardly traced to either L1 influence or target language input, thus raising questions as to whence they emerge in the first place.

This project draws on innovative methods from the psycholinguistics literature on adult native speakers to investigate key questions concerning resumptive RCs in adult L1/L2 English: First, is resumption a licit option for relativization in the grammar and/or a means of managing cognitive load during real-time sentence processing? Second, if resumption does aid in RC processing, are facilitation effects observed only in production or also in comprehension? Finally, for L2 learners (L2ers), does performance vary as a function of English proficiency?

The participant sample comprises English native speaker controls (ENSs), L1-Korean L2ers of English (KLEs), and L1-Mandarin L2ers of English (MLEs); note that only Mandarin allows resumptive RCs in any of the syntactic environments tested. Participants were distributed across two studies: one on direct object RCs (ORCs), the other on subject RCs (SRCs). Both have three main tasks: an oral elicited production task (EPT) probing processing during RC production, a self-paced reading task (SPRT) probing processing during RC comprehension, and an offline acceptability judgment task (AJT) testing acceptability of the sentence types in the other tasks. The conditions in each task target three increasingly complex RC environments: short-distance, long-distance, \textit{wh}-island. L2ers completed two versions of the AJT---one in English, one in their L1---so that L1 transfer effects, if present, could be identified. All participants also completed a language background questionnaire, an English proficiency C-test, and a short exit survey.

L1/L2 results from the EPT and the SPRT provide clear-cut evidence that resumption eases both production and comprehension of ORCs/SRCs under processing strain. For L1 adult processing, this finding contests claims in the literature (e.g.,~\cite{ferreira.swets2005}) that resumption helps only with language production. For L2 research, clear facilitation effects for resumption emerged in both sets of the processing data even when L2ers who consistently accepted resumptive RCs in the English AJT had been removed; this finding that L2 resumption is in many cases a purely processing-based phenomenon rather than a representational one casts doubt on a traditional assumption in the L2 literature (e.g.,~\cite{hyltenstam1984}) that is often implicitly held to this day.

While a sizable portion of L2ers did consistently accept English resumptive RCs, proficiency effects also arose in the AJT data: Individuals with higher C-test scores were more likely both to reject resumptive RCs---at least in the short- and long-distance environments, where English allows only gapped RCs. These results indicate that with rising English proficiency, it is possible for L2ers to become sensitive to the English prohibition on resumption in RCs.

\newpage

%%%%%%%%%%%%%%%%%%%%%%%%%%%%%%%%%%%%%%%%%%%%%%%%%%%%%%%%%%
% # Table of Contents
%%%%%%%%%%%%%%%%%%%%%%%%%%%%%%%%%%%%%%%%%%%%%%%%%%%%%%%%%%

\setstretch{1}
\renewcommand{\contentsname}{\vspace{-27pt} \begin{center} \normalsize \textbf{TABLE OF CONTENTS} \end{center} \vspace{-1.5cm}}
\renewcommand\thepage{\romannumeral\numexpr\value{page}\relax}
\phantomsection
\addcontentsline{toc}{chapter}{Table of Contents}
\tableofcontents
\newpage

%%%%%%%%%%%%%%%%%%%%%%%%%%%%%%%%%%%%%%%%%%%%%%%%%%%%%%%%%%
% # List of Tables
%%%%%%%%%%%%%%%%%%%%%%%%%%%%%%%%%%%%%%%%%%%%%%%%%%%%%%%%%%

\renewcommand{\listtablename}{\vspace{-80pt} \begin{center} \normalsize \textbf{LIST OF TABLES} \end{center} \vspace{-1.5cm}}
\phantomsection
\listoftables
\addcontentsline{toc}{chapter}{List of Tables}
\newpage

%%%%%%%%%%%%%%%%%%%%%%%%%%%%%%%%%%%%%%%%%%%%%%%%%%%%%%%%%%
% # List of Figures
%%%%%%%%%%%%%%%%%%%%%%%%%%%%%%%%%%%%%%%%%%%%%%%%%%%%%%%%%%

\renewcommand{\listfigurename}{\vspace{-77pt} \begin{center} \normalsize \textbf{LIST OF FIGURES} \end{center} \vspace{-1.5cm}}
\phantomsection
\listoffigures
\addcontentsline{toc}{chapter}{List of Figures}

%%%%%%%%%%%%%%%%%%%%%%%%%%%%%%%%%%%%%%%%%%%%%%%%%%%%%%%%%%
% # List of Abbreviations
%%%%%%%%%%%%%%%%%%%%%%%%%%%%%%%%%%%%%%%%%%%%%%%%%%%%%%%%%%

\chapter*{List of Abbreviations}
\addcontentsline{toc}{chapter}{List of Abbreviations}
\renewcommand\thepage{\romannumeral\numexpr\value{page}\relax}
\setstretch{1.5}

\noindent This dissertation makes use of the following abbreviations.

\begin{multicols}{2}
\begin{tabbing}
1 \hspace{1.5cm} \= 1st person \\
3 \> 3rd person \\
ACC \> Accusative \\
ADN \> Adnominal \\
AJT \> Acceptability judgment task \\
C \> Complementizer \\
\textit{CI} \> Confidence interval \\
CL \> Classifier \\
COMP \> Complementizer \\
CP \> Complementizer phrase \\
D \> Determiner \\
DEC \> Declarative \\
DET \> Determiner \\
DIST \> Distal demonstrative \\
DP \> Determiner phrase \\
EE \> ENS in the English AJT \\
ENS \> English native speaker \\
EPT \> Elicited production task \\
FEM \> Feminine \\
GEN \> Genitive \\
IL \> Interlanguage \\
INS \> Instrumental \\
KE \> KLE in the English AJT \\
KK \> KLE in the Korean AJT \\
KLE \> L1-Korean L2er of English \\
L1 \> First language \\
L2 \> Second language \\
L2er \> L2 learner \\
LOC \> Locative \\
M \> Masculine \\
MASC \> Masculine \\
ME \> MLE in the English AJT \\
MLE \> L1-Mandarin L2er of English \\
MM \> MLE in the Mandarin AJT \\
$n$ \> Sample size \\
N \> Noun \\
NOM \> Nominative \\
NP \> Noun phrase \\
O \> Object \\
ORC \> Direct object relative clause \\
$p$ \> $p$-value \\
PL \> Plural \\
PRED \> Predicate position \\
PRES \> Present \\
PROX \> Proximal demonstrative \\
PST \> Past \\
RC \> Relative clause \\
REL \> Relativizer \\
RM \> Relative marker \\
RP \> Resumptive pronoun \\
S \> Singular / Subject \\
\textit{SD} \> Standard deviation \\
\textit{SE} \> Standard error \\
SG \> Singular \\
SPRT \> Self-paced reading task \\
SRC \> Subject relative clause \\
T \> Tense \\
$t$ \> $t$-statistic \\
TL \> Target language \\
TOP \> Topic \\
TP \> Tense phrase \\
V \> Verb \\
VP \> Verb phrase \\
WH \>\textit{Wh}-feature \\
$z$ \> $z$-statistic \\
$\alpha$ \> First-argument position \\
$\beta$ \> Second-argument position \\
$\hat{\beta}$ \> Parameter estimate \\
\end{tabbing}
\end{multicols}

\noindent In linguistic examples with interlinear glossing, personal pronouns are represented with a number followed by one or more letters. For instance, ``1S'' marks a 1st person singular pronoun and ``3MS'' marks a 3rd person masculine singular pronoun. Complementizers are abbreviated to ``COMP'' in interlinear glosses and to ``C'' in syntactic trees. Hebrew Academy romanization is used to transcribe Hebrew sentences, Yale romanization is used to transcribe Korean sentences, and pinyin romanization is used to transcribe Mandarin sentences. Ungrammatical sentences are marked with asterisks.